% !TeX root = ../main.tex
\section{Validation and Sanity Checks}\label{sec:validation}

The microbenchmarks have no reference output to diff, so \code{scripts/validate.sh} (also
\code{make test}) gates them against physical invariants of any sane hierarchy. \code{run.sh} executes it immediately
before every experiment and stores the output as \code{validate.txt} with that dataset. The checks are:
(i)~\textbf{monotone levels:} random-chase latency grows strictly, L1 $<$ L2 $<$ L3 $<$ DRAM
(\cref{tab:zero-queue}); (ii)~\textbf{prefetcher:} sequential DRAM ``latency'' is more than $2\times$ below
random, so the prefetcher is engaged for sequential access and defeated by the randomized chase;
(iii)~\textbf{bandwidth ceiling:} L1-resident read bandwidth exceeds the DRAM-resident rate, and ten threads
exceed one; (iv)~\textbf{loaded $>$ unloaded:} the \code{mem_loaded} probe is slower under four streaming
cores than idle; (v)~\textbf{frequency plausibility:} the add-chain probe lands in
\qtyrange{1}{6}{\giga\hertz} (coarse only; conversions use the fixed clock). Reported but not gating:
counters open and return values, and a \qty{2}{\mebi\byte}-page request is granted (\code{AnonHugePages}
$>0$) while the \qty{4}{\kibi\byte} policy yields none.

\paragraph{External cross-check: Intel MLC.} \code{run.sh mlc} runs MLC v3.13 (\code{-e -r}) with its
latency thread on CPU~2 and load threads on the other nine cores. \Cref{tab:zero-queue} compares its idle
latency with ours and \cref{sec:intensity} its loaded-latency sweep with our knee.

\paragraph{Additional internal checks.} Every zero-queue size is measured twice: an LCG ring timed with
\code{CLOCK_MONOTONIC_RAW} and an xorshift64* ring timed with \code{rdtscp}; on the DRAM plateau they differ by
\val{zq_DRAM_xdiff}\% (\code{validate.sh} fails above 15\%). The separately written pointer ring of
\code{mem_loaded} gives \valqty{int_w0_ns}{\nano\second} unloaded, against \valqty{zq_DRAM256_ns}{\nano\second}
from \code{mem_latency} at \qty{256}{\mebi\byte}. The L1 latency of \valnum{zq_L1_cyc}~cycles matches this
core's published 4--5 cycle load-to-use latency, anchoring the ns$\to$cycle conversion.
